% META: {"id": "1NSI-ARCHITECTURE-OS-RE-C5", "chapitre": "1NSI-ARCHITECTURE-OS", "type_objet": "remediation", "capacites": ["P-ARCH-03C"], "status": "generated"}

%---------------------------------------------------------------
% FICHE DE REMEDIATION — C5 : confondre les 3 blocs de permissions (proprietaire/groupe/autres)
%---------------------------------------------------------------

\begin{exercice}{1NSI-ARCHOS-RE-C5-EX1}{1}{10}
Un élève lit la permission \lstinline{"rwxrwxr--"} et affirme que « le bloc du milieu
(\lstinline{rwx}) concerne les \emph{autres} utilisateurs », alors qu'il concerne en
réalité le \emph{groupe}.
\begin{enumerate}
  \item Rappeler l'ordre des trois catégories dans une chaîne de permissions.
  \item En utilisant \lstinline{droit_accorde} du cours, déterminer si le \textbf{groupe} peut écrire dans ce fichier, puis si les \textbf{autres} utilisateurs peuvent y écrire.
  \item Ces deux réponses sont-elles identiques ? Qu'est-ce que cela montre sur l'erreur de l'élève ?
\end{enumerate}
\end{exercice}

% BEGIN-VERIFY
% def droit_accorde(permissions, categorie, action):
%     index = {"proprietaire": 0, "groupe": 3, "autres": 6}
%     decalage = {"lecture": 0, "ecriture": 1, "execution": 2}
%     position = index[categorie] + decalage[action]
%     return permissions[position] != "-"
% perm = "rwxrwxr--"
% assert droit_accorde(perm, "groupe", "ecriture") is True
% assert droit_accorde(perm, "autres", "ecriture") is False
% END-VERIFY

\begin{corrige}{1NSI-ARCHOS-RE-C5-EX1}
\textbf{1.} L'ordre est toujours : \textbf{propriétaire} (caractères 1 à 3),
\textbf{groupe} (caractères 4 à 6), \textbf{autres} (caractères 7 à 9).

\textbf{2.} Pour \lstinline{"rwxrwxr--"} : le bloc \textbf{groupe} est
\lstinline{rwx} (caractères 4-6), donc le groupe \textbf{peut} écrire
(\lstinline{droit_accorde(perm, "groupe", "ecriture")} vaut \lstinline{True}). Le bloc
\textbf{autres} est \lstinline{r--} (caractères 7-9), donc les autres \textbf{ne peuvent
pas} écrire (\lstinline{False}).

\textbf{3.} Non, ces deux réponses sont \textbf{différentes} (\lstinline{True} contre
\lstinline{False}) : cela confirme que le bloc du milieu et le dernier bloc concernent des
catégories distinctes, avec des droits potentiellement différents. L'élève avait tort de
considérer le bloc du milieu comme celui des \og autres \fg{} : c'est celui du
\textbf{groupe}.
\end{corrige}
